\anonsection{Заключение}
\label{sec:7_zakluch} \index{7_zakluch}

\singlespace

В результате выполнения работы получены следующие основные результаты:
\begin{enumerate}
	\item Разработан жадный алгоритм списочного планирования для многопроцессорной системы с ограничением на суммарный объём памяти. Реализованы две эвристики выбора работы: EFT (минимизация времени завершения работы) и STS (учёт затрат памяти на работу).
	\item Выполнено сведение задачи к задаче целочисленного линейного программирования (ЦЛП) с использованием переменных, моделирующих пересечение интервалов активности буферов. Реализовано построение расписания при помощи ЦЛП.
	\item Проведено экспериментальное исследование на трёх классах графов (слоистые, случайные, треугольные) размером от 10 до 990 вершин. Показано, что:
	\begin{itemize}
		\item ЦЛП находит оптимальное решение для всех малых графов с числом вершин не более 21, причем наибольшее преимущество у ЦЛП над лучшим из ЖА составляет более 32\%.
		\item При числе вершин больше 16 и малом числе доступных процессоров время работы ЦЛП приближается к лимиту 3600~с; для больших графов ЦЛП не применимо.
		\item Жадные алгоритмы работают на порядки быстрее, чем ЦЛП (миллисекунды–секунды). Среднее отклонение длительности ЖА от длительности ЦЛП составляет $9$–$13\%$, максимальное – до $35\%$.
		\item Сравнение эвристик EFT и STS показало, что STS лучше на треугольных графах и при больших ресурсах (числе процессоров, объеме памяти), но уступает EFT на случайных графах с малым числом процессоров. Рекомендуется запускать обе эвристики и выбирать лучшее расписание.
	\end{itemize}
\end{enumerate}

Направления дальнейших исследований включают:
\begin{itemize}
	\item поддержку целевых систем с процессорами различной производительности (длительность работы зависит от процессора);
	\item применение методов машинного обучения для выбора эвристики или предсказания «хороших» упорядочений переменных для ЦЛП;
	\item разработку гибридных алгоритмов, комбинирующих жадные эвристики и точные методы для графов среднего размера (50–200 вершин).
\end{itemize}